\section{Discussion}
\label{sec:discussion}

\subsection{Scope of patch positivity}

Patch positivity concerns the conditional contribution of the dual updates
within one patch, not the sign of each update separately. Consecutive updates
in the unaveraged Feynman--Kac dual may have different signs. After
conditioning on the successful-interaction skeleton, the remaining local marks
are averaged inside each patch to give~$C(P)$. The coefficient criterion
determines exactly when these patch contributions are nonnegative.

The condition is neither configuration monotonicity nor positivity of every
coefficient in the signed dual. When~$\mb p^\star=\mb 0$, the centered
monomials are ordinary monomials and~$\preceq_*$ is the upper-orthant, or
monomial-moment, order. Stochastic order implies this order, but the converse
fails. For the contact process, monotonicity in configuration order therefore
gives stronger comparisons for all increasing observables. The centered-moment
comparison also applies to FA-1f and BABP, whose dynamics do not preserve the
configuration order.

The profile~$\mb p^\star$ is the smallest terminal reference profile for which
every end-patch contribution is nonnegative. It records where the sign argument
applies. The lower profile
$\mb p^-=(\max\cb{2p_i^\star-1,0})_{i\in\Lambda}$ arises only from the affine
extension of~$\mathcal M_*$ to~$\mathcal M_-$. Neither profile is asserted to
mark a change in the qualitative dynamics.

\subsection{Limitations of the convergence argument}

The restriction to~$\mathcal M_-$ comes from the centered expansion of the end
factors. Outside this class some centered terms are negative, so the proof has
no termwise nonnegative comparison. It therefore gives no conclusion about
convergence outside~$\mathcal M_-$. Such an extension would have to control
cancellations between terms of different signs, uniformly over the successful
skeleton. When~$\mb p^\star\leq\mbs{\tfrac12}$, every initial law belongs
to~$\mathcal M_{-,1}$ and this restriction disappears.

The uniform pure-death component is used twice. Along a backward chain of
outgoing patches it supplies a survival factor for realizations with a
successful interaction after the cut time. When there is no such interaction,
it contracts the centered slopes of the end patches and removes their
dependence on the terminal configuration. The proof would extend to a hard
model if model-specific estimates replaced these two bounds.

The full-patch formula can sometimes identify the invariant measure. If every
infinite outgoing-to-end patch has zero contribution, every realization with an
ordinary successful interaction has zero limiting weight. Indeed, on
$\cb{\abs{\mathcal P}<\infty}$, such a realization has a last
successful interaction, and the source patch beginning there is an infinite
outgoing-to-end patch. The remaining term
is the independent spin system with~$0\to1$ and~$1\to0$ rates
$c_i^0(\vn)$ and~$c_i^1(\vn)$, whose stationary probability of state~$1$ is

\begin{equation*}
\frac{c_i^0(\vn)}{c_i^0(\vn)+c_i^1(\vn)}.
\end{equation*}

This recovers~$\mu_p$ for the equilibrium-refresh perturbations of FA-1f and
BABP. Without the vanishing condition, the representation still gives the
patchwise invariant-moment comparison in
Corollary~\ref{cor:pure-death-invariant-comparison}.

Two problems remain. The first is convergence from arbitrary initial laws when
$\mb p^\star\nleq\mbs{\tfrac12}$. The second is convergence from configurations
with finitely many sites in state~$0$ without a uniform pure-death component,
including FA-1f from finitely many vacancies and BABP from a finite nonempty
particle set. The latter requires model-specific control because these point
masses are not bracketed by useful nondegenerate product laws under
$\preceq_*$ and the hard dynamics has no uniform pure-death component.
